\subsection{弱完备空间}

命题 9. --- 设 F、G 为两个处于分离对偶的向量空间。若 \(\hat{\mathbf{F}}\) 是空间 F 关于拓扑 \(\sigma(\mathrm{F},\mathrm{G})\) 的完备化，且我们考虑典范单射 \(j:\mathrm{F}\to\mathrm{G}^{*}\)，其中 \(\mathrm{G}^{*}\) 赋拓扑 \(\sigma(\mathrm{G}^{*},\mathrm{G})\)，则连续延拓 \(\hat{j}:\hat{\mathrm{F}}\to\mathrm{G}^{*}\)（延拓自 \(j\)）是拓扑向量空间的一个同构。

因为，我们看到 \(G^{*}\) 赋以 \(\sigma(G^{*}, G)\) 后是 Hausdorff 的且完备（II, p.~51, 推论 2）；若借助 j 把 F 等同于 \(G^{*}\) 的一个向量子空间，则 \(\sigma(G^{*}, G)\) 在 F 上诱导的拓扑是 \(\sigma(F, G)\)，且 F 在 \(G^{*}\) 中关于拓扑 \(\sigma(G^{*}, G)\) 稠密（II, p.~43, 推论 4）；命题由此得证。

于是，关于某个弱拓扑完备的向量空间，就是任意的向量空间 G 的对偶 G* 赋以 \(\sigma(\mathbf{G}^{*}, \mathbf{G})\)；由 II, p.~51 的推论 2，它们（拓扑上）同构于实直线的乘积 \(R^{I}\)。为使措辞简化，我们将称它们为「直线积」（关于这类空间的内蕴刻画，见 II, p.~85, 习题 13 及 II, p.~81, 习题 1）。

我们注意到，在 \(G^{*}\) 上，拓扑 \(\sigma(G^{*}, G)\) 是 Hausdorff 的弱拓扑中最小的；因为，一个粗于 \(\sigma(G^{*}, G)\) 的弱拓扑必然形如 \(\sigma(G^{*}, H)\)，其中 \(H \subset G\)（II, p.~43, 推论 3）；但若 \(H \neq G\)，则存在一个线性型 \(x^{*} \in \mathbf{G}^{*}\)，它非零且与 \(\mathbf{H}\) 正交（A, II, § 7.3, 命题 8），因此 \(\sigma(\mathbf{G}^{*}, \mathbf{H})\) 不是 Hausdorff 的。

由这个注我们推出：若 F、G 是两个向量空间，则线性双射 \(u:G^{*}\to F^{*}\) 只要关于拓扑 \(\sigma(G^{*},G)\) 与 \(\sigma(F^{*},F)\) 连续，就必然是双连续的。

命题 10. --- 设 G 为实向量空间，F = G* 是它的对偶，赋拓扑 \(\sigma(\mathbf{G}^*, \mathbf{G})\)。

\begin{enumerate}
\def\labelenumi{(\roman{enumi})}
\item
  映射 \(\mathbf{V} \mapsto \mathbf{V}^{\circ}\) 是 \(\mathbf{G}\) 的向量子空间全体到 \(\mathbf{F}\) 的闭向量子空间全体上的双射。
\item
  \(\mathbf{F}\) 的每个闭向量子空间都是直线积且具有拓扑补。
\end{enumerate}

由双极定理（II, p.~45, 推论 2），\(\mathrm{V} \mapsto \mathrm{V}^{\circ}\) 是 G 中关于 \(\sigma(\mathbf{G}, \mathbf{G}^{*})\) 为闭的向量子空间 V 全体到 F 的闭向量子空间全体上的双射。但按定义，G 上每个线性型都关于 \(\sigma(\mathbf{G}, \mathbf{G}^{*})\) 连续，因此 G 中每个向量子空间都是闭的，因为它由一个方程组 \(\langle y, y_{\lambda}^{*} \rangle = 0\)（其中 \(y_{\lambda}^{*} \in \mathbf{G}^{*}\)）定义；这就证明了 (i)。

现在设 W 是 F 的一个闭子空间；于是有 \(W = V^{\circ}\)，其中 \(V = W^{\circ}\) 在 G 中。设 \(V'\) 是 V 在 G 中的一个补空间。我们知道 \(F = G^{*}\) 可典范地等同于 \(V^{*} \oplus V'^{*}\)，且 \(V'^{*}\) 等同于 \(V^{\circ} = W\)（A, II, § 2.6, 命题 10 的推论）；此外（II, p.~50, 命题 8），拓扑 \(\sigma(G^{*}, G)\) 可等同于拓扑 \(\sigma(V^{*}, V)\) 与 \(\sigma(V'^{*}, V')\) 的乘积；这就证明了断言 (ii)。

尽管对拓扑 \(\sigma(\mathbf{G},\mathbf{G}^{*})\) 而言 \(\mathbf{G}\) 的每个向量子空间都是闭的，我们仍注意到，若 \(\mathbf{G}\) 是无限维的，则拓扑 \(\sigma(\mathbf{G},\mathbf{G}^{*})\) 并不是 \(\mathbf{G}\) 上最细的局部凸拓扑，因为 \(\sigma(\mathbf{G},\mathbf{G}^{*})\) 的每个零邻域都包含一个无限维向量子空间：然而它是 \(\mathbf{G}\) 上弱拓扑中最细的（II, p.~43, 推论 3）。

